% Fiche de préparation de la visite à la BnF (#37) : les faits à vérifier,
% point par point, devant les boîtes de NAF 29093. À imprimer, à remplir au
% crayon (les salles des manuscrits n'admettent pas les stylos).
\documentclass[10pt,a4paper]{article}
\usepackage{fontspec}
\IfFontExistsTF{Charter}{\setmainfont{Charter}}{\setmainfont{XCharter}}
\usepackage{polyglossia}\setmainlanguage{french}
\usepackage[a4paper,margin=16mm,top=14mm]{geometry}
\usepackage{tabularx,array,booktabs,enumitem,xcolor,amssymb}
\usepackage[hidelinks]{hyperref}
\setlength{\parindent}{0pt}\setlength{\parskip}{3pt}
\setlist{nosep,leftmargin=*}
\newcommand{\cb}{$\square$}
\newcommand{\ligne}{\rule{0pt}{4.2ex}}
\newcommand{\q}[1]{«\,#1\,»}
\newcolumntype{Y}{>{\raggedright\arraybackslash}X}
% Une ligne de vérification : le fait annoncé, puis conforme / différent / à noter.
\newcommand{\fait}[1]{#1 & \cb & \cb & \\ \hline}
\newenvironment{faits}{\renewcommand{\arraystretch}{1.35}\par\noindent\tabularx{\linewidth}{|Y|c|c|p{0.32\linewidth}|}\hline
\textbf{Fait annoncé (notice BnF, issue \#37)} & \textbf{Oui} & \textbf{Non} & \textbf{Ce que montre la boîte}\\ \hline}{\endtabularx\par\smallskip}
\newcommand{\boite}[2]{\subsection*{#1}{\small\color{gray!60!black}#2}\par}

\begin{document}
{\Large\bfseries Fiche de visite --- BnF, NAF 29093}\par
{\small Fonds Alexandre Grothendieck, «\,Œuvres\,» · Département des Manuscrits (Richelieu) · demande à \texttt{manuscrits@bnf.fr} · issue GitHub \#37 · mise à jour le 3 octobre 2026}\par
\medskip
\textbf{Date de la séance :} \rule{2.8cm}{0.4pt}\quad \textbf{Salle :} \rule{2.5cm}{0.4pt}\quad \textbf{Boîtes servies :} \hrulefill

\section*{Avant de partir}
\begin{tabularx}{\linewidth}{@{}p{0.5\linewidth}Y@{}}
\cb\ Réponse de \texttt{manuscrits@bnf.fr} reçue (accord, conditions)  & \cb\ Pièce d'identité avec photo \\
\cb\ Pass Recherche : entretien à l'Espace conseil prévu (guichets jusqu'à 18 h) & \cb\ Kbis de moins de 3 mois, lettre à en-tête de la SASU \\
\cb\ Liste des documents à consulter (ci-dessous), imprimée & \cb\ Crayons à papier seulement \\
\cb\ Photographies : ce qui est permis, demandé par écrit & \cb\ Index et timeline du site, hors ligne
\end{tabularx}

\textbf{Liste à remettre :} NAF 29093 (3), (43), (38--39), (4--6), (40), (42), dans cet ordre de priorité.

\section*{1. Le fonds dans son ensemble}
\begin{faits}
\fait{Cote NAF 29093, intitulé \q{Fonds Alexandre Grothendieck. Œuvres.}}
\fait{43 boîtes, de 1987 à 1999, en français}
\fait{Entré en 2023, par \q{legs de son auteur complété par un don de ses enfants}}
\fait{Non numérisé ; communication soumise à une \q{procédure spécifique}}
\fait{Un seul volume présenté comme mathématique : (3)}
\end{faits}
Conditions de communication obtenues sur place : \hrulefill\par\ligne\rule{\linewidth}{0.4pt}

\section*{2. Les boîtes à demander, dans l'ordre}

\boite{(3) \emph{Géométrie élémentaire schématique}}{Priorité 1. Le test direct de la question de l'article RHM : les mathématiques continuent-elles après Montpellier, et dans quelle langue ?}
\begin{faits}
\fait{Daté d'août à octobre 1992}
\fait{\q{Recueil de résultats mathématiques auxiliaires}, en lien avec une \q{variété riemanno-lorentzienne}}
\fait{\q{Prolongeant les travaux entamés \ldots\ dans ses \emph{Éléments de géométrie algébrique}} : langage des schémas, renvois à EGA}
\fait{Manuscrit autographe (de sa main)}
\fait{Feuillets datés au jour, comme en 1983--1986}
\end{faits}
Notation employée (schémas, faisceaux, topos, catégories, \q{formes}) : \hrulefill\par\ligne\rule{\linewidth}{0.4pt}

\boite{(43) Dossiers de notes et \q{gribouillis}}{Priorité 2. La boîte commence là où s'arrête Montpellier (les \emph{Dérivateurs}, 157, 1990--1991).}
\begin{faits}
\fait{Daté d'août 1991 à février 1997}
\fait{Notes préparatoires, brouillons de sections, \q{suites de schémas et diagrammes}}
\fait{Un fil de Montpellier s'y poursuit : dérivateurs (157), géométrie des formes (156), pseudo-droites (154)}
\fait{Les motifs y reviennent (absents du corpus tardif de Montpellier)}
\end{faits}
Premier et dernier feuillet datés : \hrulefill

\boite{(38--39) \q{Tables des matières} et \q{notes-résumés}}{Priorité 3. Sa propre carte du \emph{Problème du Mal} et des \emph{Réflexions} : mesurer le reste sans lire des dizaines de milliers de feuillets.}
\begin{faits}
\fait{Tables des matières de l'ensemble des \emph{Réflexions sur la Vie et le Cosmos}}
\fait{Des chapitres ou sections y portent sur les mathématiques}
\fait{Les \emph{Réflexions mathématiques} (plan de 1983) y sont annoncées ou mentionnées}
\end{faits}
Chapitres mathématiques repérés (boîte, titre, feuillets) : \par\ligne\rule{\linewidth}{0.4pt}\par\ligne\rule{\linewidth}{0.4pt}

\boite{(4--6) \emph{Structures de la Psyché} \ ·\ (7) \emph{Psyché et Structure}}{Priorité 4. Un texte \q{rédigé en langage formel}.}
\begin{faits}
\fait{Daté d'octobre 1992 à avril 1993}
\fait{\q{Rédigée en langage formel} : définitions, énoncés, démonstrations}
\fait{Emprunte aux catégories, aux topos ou aux \q{formes} de 156}
\fait{(7) est le texte littéraire qui l'accompagne}
\end{faits}

\boite{(40) et (42) \q{Gribouillis} I et II}{Priorité 5. Les diagrammes, à comparer aux dessins tardifs de Montpellier (154, 156 ; issue \#40).}
\begin{faits}
\fait{Datés de 1997 à 1999}
\fait{Des diagrammes ou figures mathématiques (graphes, cartes, pavages)}
\fait{Ressemblance avec les dessins de 154 ou 156}
\end{faits}

\boite{Plus tard, seulement s'il y a des mathématiques}{}
\begin{faits}
\fait{(14--15, 34) \emph{Esprit et Matière}, 1992--1996 : physique, des équations de Maxwell à la relativité générale}
\fait{(1--2) \emph{La Clef des Songes}, dactylographie corrigée, avril 1987 -- avril 1988 ; (2) = notes du chapitre VII, \q{Mutants}}
\fait{(8--13, 16--33, 35--39) \emph{Le Problème du Mal}, 1993--1998, \q{plusieurs dizaines de milliers de feuillets}}
\fait{(41) \emph{Réflexions sur la Création}, 1998--1999}
\end{faits}

\section*{3. Grille par boîte}
\renewcommand{\arraystretch}{1.9}
\begin{tabularx}{\linewidth}{|p{1.2cm}|p{1.6cm}|p{2.2cm}|c|Y|Y|p{1.4cm}|}\hline
\textbf{Boîte} & \textbf{Étendue} & \textbf{Dates (de sa main ?)} & \textbf{\% math.} & \textbf{Sujets (index du site)} & \textbf{Noms cités} & \textbf{Photos ?} \\ \hline
(3) &&&&&& \\ \hline (43) &&&&&& \\ \hline (38) &&&&&& \\ \hline (39) &&&&&& \\ \hline
(4) &&&&&& \\ \hline (5) &&&&&& \\ \hline (6) &&&&&& \\ \hline (40) &&&&&& \\ \hline (42) &&&&&& \\ \hline
\end{tabularx}

\section*{4. Les questions de l'article RHM}
\begin{faits}
\fait{Il revient aux motifs après 1991 (RHM §\,6)}
\fait{Les fils de Montpellier se poursuivent (157, 156, 154)}
\fait{Il date encore ses feuillets au jour}
\fait{Des lettres et des interlocuteurs mathématiques après 1991}
\fait{Les \emph{Réflexions mathématiques} annoncées en 1983 existent sous une forme}
\fait{Les correspondants de Montpellier reviennent : Breen (lettres de 1975, cote 134-2), Quillen (19.2.1983), Deligne, Illusie}
\end{faits}

\section*{Où en est Montpellier (3 octobre 2026)}
{\small 154 cotes commencées, 7\,783 pages transcrites en 538 lots ; dix cotes non éditées restent à transcrire (\#43) ; 81 cotes ont une lecture modernisée. Le fil de comparaison le plus serré est \emph{À la poursuite des champs} (134), les \emph{Dérivateurs} (157, 158) et la géométrie des formes (156), dont les dernières pages datées sont de 1986 à 1991.}

\section*{Après la séance}
\cb\ Compte rendu écrit, hors du dépôt public si la BnF le demande \quad \cb\ Limites de l'article RHM mises à jour \quad \cb\ \#30 réexaminée \quad \cb\ Rien de publié sans autorisation

\end{document}
